\documentclass[11pt, letterpaper]{article}

\input{/home/hyperion/.config/LatexHub/preambles/base.tex}
\input{/home/hyperion/.config/LatexHub/preambles/tikz.tex}
\input{/home/hyperion/.config/LatexHub/preambles/diagrams.tex}
\input{/home/hyperion/.config/LatexHub/preambles/macros-cat-theory.tex}
\input{/home/hyperion/.config/LatexHub/preambles/macros-algebra.tex}
\input{/home/hyperion/.config/LatexHub/preambles/refs-basic.tex}
\input{/home/hyperion/.config/LatexHub/preambles/basic-theorems-section.tex}

\theoremstyle{theorem}
\newtheorem{problem}{Problem}

% TODO
% 1.7 pain
% 2.3
% 2.6 smoothnesss
% 2.9 smoothness
% 2.10 (2, extend psi)
% get rid of categoricals
% ok quotient maps arent open in genearl 
\usepackage{titlepageBU}
\title{Homework 1}
\begin{document}
\makeproblem

\begin{center}
	Due to some poor time management I did not complete all the problems. Obviously won't happen again; I just misjudged my time committments this semester. I've marked problems where I failed to complete them for the grader's convenience.
\end{center}


\section{Problem 1.1}
\begin{problem}\label{1}
	Let $X\subseteq \mathbb{R} ^2$ be the set of all points $(x,y)$ with $y=\pm 1$, and let
	\[
		M\coloneqq X / \left\langle (x,-1)\sim (x,1) \text{ when } x\neq 0 \right\rangle 
	.\]
	Show that $M$ is locally Euclidean and second countable, but not Hausdorff.
\end{problem}

\begin{definition}[Lee A.1.3]\label{1.1}
	Let $\left\{ X_\alpha  \right\}_{\alpha \in A} $ be a collection of spaces. The disjoint union is topologized by taking a basis $\mathcal{B} $ to consist of the images of each open set in each $X_\alpha $ under the canonical inclusions.
\end{definition}


\begin{proposition}\label{1.2}
	Let $X,Y$ be spaces with bases $\mathcal{B} _X,\mathcal{B} _Y$, respectively. Then $X\amalg Y$ has a basis given by $\mathcal{B}'\coloneqq \left\{ U\subseteq X\amalg Y\mid U\in \mathcal{B} _X \text{ or } U\in \mathcal{B} _Y \right\} $.
\end{proposition}
\begin{proof}
	It suffices to show the basis $\mathcal{B} $ given in \cref{1.1} is generated by $\mathcal{B}'$. Let $B\in \mathcal{B} $, and without loss of generality let $B\subseteq X$. By hypothesis, there exists a collection $\left\{ B_i \right\}_{i\in I} \subseteq \mathcal{B} _X\subseteq \mathcal{B} '$ such that
	\[
		B = \bigcup_{i\in I} B_i
	.\]
\end{proof}

\begin{remark}\label{1.3}
	Functions preserve unions, but for $f : X \to Y$ and $U,V\subseteq X$, we only have $f(X \cap Y)\subseteq f(X)\cap f(Y)$. Preimages preserve both unions and intersections. Injections preserve intersections. This is a well-known fact of set theory that I do not feel like proving for the third semester in a row.
\end{remark}


\begin{lemma}\label{1.4}
	Second countability is a topological invariant. More generally, open surjections preserve second countability.
\end{lemma}
\begin{proof}
	Let $f : X \to  Y$ surjective with $X$ second countable. Let $\mathcal{B} $ be a countable basis for $X$. Define
	\[
		f[\mathcal{B} ]\coloneqq \left\{ f(B)\mid B\in \mathcal{B}  \right\} 
	.\]
	Since $f$ is open, $f(U)$ is open for all $U\subseteq X$ open. Clearly the cardinalities are such that $\left| f[\mathcal{B} ] \right| \leq \left| \mathcal{B}  \right| $, so it suffices to show $f[\mathcal{B} ]$ is a basis for $Y$. Let $y\in Y$. Then there exists at least one $x\in X$ with $f(x)=y$ by hypothesis. Since $\mathcal{B} $ is a basis for $X$, there exists an open set $x\in U\subseteq X$, and thus $y\in f(U)\in f[\mathcal{B} ]$. This gives us an open cover of $Y$.

	Let $f(B_1),f(B_2)\in f[\mathcal{B} ]$, and let $y=f(x)\in f(B_1)\cap f(B_2)$. By \cref{1.3}, we have
	\[
		x\in f^{-1} (f(B_1)\cap f(B_2))=f^{-1} f(B_1)\cap f^{-1} f(B_2)
	.\]
	Since $f$ is continuous, $f^{-1} f(B_1),f^{-1} f(B_2)$ are open in $X$. $\mathcal{B} $ is a basis for $X$, there exists $B_3\in \mathcal{B} $ with $x\in B_3 \subseteq f^{-1} f(B_1) \cap f^{-1} f(B_2)$. Then
	\begin{align*}
		f(x)=y\in f(B_3)&\subseteq f(f^{-1} f(B_1)\cap f^{-1} f(B_2))\\
				&=ff^{-1} (f(B_1)\cap f(B_2))\\
				&=f(B_1)\cap f(B_2)
	,\end{align*}
	and since $f(B_3)\in f[\mathcal{B} ]$, we have that $f[\mathcal{B} ]$ is a basis.
\end{proof}


\begin{proof}[Proof of \cref{1}]
	It should be extremely obvious that $X \cong \mathbb{R} \amalg \mathbb{R} $. Since $\mathbb{R} $ is second countable, let $\mathcal{B} $ be a countable basis for $\mathbb{R} $. By \cref{1.2}, $\mathcal{B} \amalg \mathcal{B} $ (minor abuse of notation) is a basis for $\mathbb{R} \amalg \mathbb{R} $. By basic set theory, $\mathcal{B} \amalg \mathcal{B} $ is countable. By \cref{1.4}, $X$ is second countable. By \cref{1.4}, it suffices for second-countability to show the quotient map $\pi $ is open. We delegate this to after the proof that it is locally Euclidean.

	We now show that $M$ is locally Euclidean. Of course, if $x\neq 0$, then the statement is fairly trivial: let $r=d(x,0)$, and note that the ball $B_r(x)$ considered as a subset of $M$ in the obvious way is open (its preimage is $B_r(x)\amalg B_r(x)$). Of course we have $M\supseteq B_r(x)\cong B_r(x)\subseteq \mathbb{R} $. A similar proof also holds for $x=0$, where $0$ can be an element of either copy of $\mathbb{R} $. Write $0_1$ for the image of $(0,1)\in X $ in $M$, and similarly write $0_{-1} $ for the image of $(0,-1)\in X$. Consider the set $B_\varepsilon (0_1)$ defined by
	\[
		B_\varepsilon (0_1)=\left\{ x\in M \mid 0<d(x,0_1)<\varepsilon  \right\} \cup \left\{ 0_1 \right\} ,
	\]
	where the metric $d(-,-)$ is defined in the obvious way: write $\pi' : M \to \mathbb{R} $ for the quotient by $0_1\sim 0_{-1} $, and $d_\mathbb{R} $ the metric on $\mathbb{R} $. Define $d= d_{\mathbb{R} } \circ \pi $. Of course, intuitively this is just the open ball about the 0 at $1$ of radius $\varepsilon $ not including the $0$ at $-1$. The preimage under $\pi : \mathbb{R} \amalg \mathbb{R}  \to M$ is very clearly just
	\[
		\pi ^{-1} (B_\varepsilon (0_1))=B_\varepsilon (0)\amalg (B_{\varepsilon /2} (\varepsilon /2)\cup B_{\varepsilon /2} (-\varepsilon /2))
	.\]
	Both sets on either side of the coproduct are obviously open, and thus the preimage is open since the ``disjoint union'' of these sets is really just the union, viewed inside a coproduct space.

	We may now use these results to streamline the proof of openness of $\pi $. Let $B_r(x)$ be an open set in one copy of $\mathbb{R} $ in the disjoint union, and note that if $0\notin B_r(x)$, we obviously have $\pi ^{-1} \pi (B_r(x))=B_r(x)\amalg B_r(x)$. This is open, and thus so is $\pi (B_r(x))$. If $0\in B_r(x)$, then the preimage is still $B_r(x)$ in the original copy of $\mathbb{R} $, but in the other copy, we have $B_r(x)\setminus \left\{ 0 \right\} $. Of course, we showed above that this is just a union of two other open balls (we centered those balls at 0 but they can be moved around and it still follows), so that set is still open. Thus, $\pi (B_r(x))$ is open, and $\pi $ is an open map. The conclusion follows by \cref{1.4}.

	To show that $M$ is not Hausdorff, consider the points $0_1,0_{-1} $. Since open sets for these are generated by images of balls under $\pi $ by \cref{1.4}, it suffices to show no images of balls containing $0_1$ and $0_{-1} $ are disjoint. Suppose $0_1\in B_r(x),0_{-1} \in B_s(y)$ are disjoint neighborhoods, where we use the ball notation in $M$ and simply imply that they do not contain both zeros. Let $-x<\varepsilon <r-x,s-y$. Since $0-x<r-x$ and $0-y<s-y$ (otherwise they wouldn't be in the neighborhoods), such a $\varepsilon $ exists. We then obviously have that $\varepsilon \in B_r(x),B_s(y)$ since $\varepsilon $ is identified under $\pi $ from both copies of $\mathbb{R} $, a contradiction. Thus, $M$ is not Hausdorff.
\end{proof}

\section{Problem 1.2}


\begin{problem}\label{2}
	Show that a disjoint union of uncountably many copies of $\mathbb{R} $ is locally Euclidean and Hausdorff, but not second countable.
\end{problem}
\begin{remark}\label{2.1}
	A standard fact of topology is that given a space $(X,\mathcal{T} )$, a set $\mathcal{B} \subseteq \mathcal{T} $ is a basis for $\mathcal{T} $ if and only if for each open $U\subseteq X$ and each point $x\in U$, there is a set $B\in \mathcal{B} $ such that $x\in B\subseteq U$.
\end{remark}

\begin{lemma}\label{2.2}
	Let $\left\{ X_\alpha  \right\}_{\alpha \in A} $ be a collection of Hausdorff spaces. Then so is $X\coloneqq \coprod_{\alpha \in A} X_\alpha $.
\end{lemma}
\begin{proof}
	Let $x,y\in X$ with $x\neq y$. Without loss of generality, let $x\in X_\alpha $ and $y\in X_\beta $ for $\alpha ,\beta \in A$. If $\alpha \neq \beta $, then we clearly have that $X_\alpha $ and $X_\beta $ are disjoint open neighborhoods. Otherwise, assume without loss of generality that $x,y\in X_\alpha $. The conclusion follows by definition of the disjoint union since $X$ is Hausdorff.
\end{proof}

\begin{proof}[Proof of \cref{2}]
	Let $I$ be a set with $\left| I \right| >\left| \mathbb{N}  \right| $, and define $X\coloneqq I\cdot \mathbb{R} $ for $\cdot $ the copower. Denote by $\mathbb{R} _i$ the copy of $\mathbb{R} $ indexed by $i\in I$. For any $x\in X$, we have $x\in\mathbb{R} _i$ for some $i$, and there is a canonical isomorphism $\mathbb{R} _i \cong \mathbb{R} $. This shows $X$ is locally Euclidean since $\mathbb{R} _i$ is open by \cref{1.1}. By \cref{2.2}, $X$ is Hausdorff.
	
	To show $X$ is not second countable, let $\mathcal{B} $ be a basis for $X$. Since $\mathbb{R} _i$ is open for all $i\in I$, by \cref{2.1} we have that for all $i$ and for all $x\in\mathbb{R} _i$, there is a set $B_i\in \mathcal{B} $ such that $x\in B_i\subseteq \mathbb{R} _i$. In particular, for each $i$ there is at least one element $B_i\in \mathcal{B} $. These sets are all distinct since $\mathbb{R} _i \cap \mathbb{R} _j$ is nonempty if and only if $i=j$. We thus have a bijection $I\cong \left\{ B_i \right\}_{i\in I} \subseteq \mathcal{B} $, and since $I$ is uncountable, $\mathcal{B} $ is at least the cardinality of $I$, and thus uncountable.
\end{proof}

\section{Problem 1.7}
\begin{problem}\label{3}
	Let $N\in S^n\subseteq \mathbb{R} ^{n+1} $ be the pole $N=(0, \ldots ,0,1)$, which we call the \emph{north pole}. Similarly, let $S\in S^n$ denote the \emph{south pole} $S=(0, \ldots ,0,-1)=-N$. Define the stereographic projection $\sigma : S^n\setminus \left\{ N \right\} \to \mathbb{R} ^n$ by
	\[
		(x^1, \ldots ,x^{n+1} ) \mapsto \frac{(x^1, \ldots ,x^n)}{1-x^{n+1} }
	.\]
	Denote by $\widetilde{\sigma } (x)=-\sigma (-x)$ for $x\in S^n\setminus \left\{ S \right\} $.
	\begin{enumerate}
		\item For any $x\in\dom \sigma $, let $(u,0)$ be the point at which the line through $N$ and $x$ intersects the linear subspace where $x^{n+1} =0$. Show $\sigma (x)=u$. Similarly, $\widetilde{\sigma } (x)$ is the point where the line from $S$ to $x$ intersects the same subspace.
		\item Show that $\sigma $ is bijective with
			\[
				\sigma ^{-1} (u^1, \ldots ,u^n) = \frac{(2u^1, \ldots ,2u^n, \left| u \right| ^2-1)}{\left| u \right| ^2+1} 
			.\]
		\item The transition map $\widetilde{\sigma } \circ \sigma ^{-1} $ (given by explicit computation) is smooth, and the two charts $\sigma ,\widetilde{\sigma } $ combine to form a smooth structure on $S^n$.
		\item This smooth structure is the same as the one given by the $\varphi _j^{\pm} $s.
	\end{enumerate}
\end{problem}
We spend some time establishing the useful result that products in the category-theoretic sense of smooth maps are smooth. This was a useful fact fairly often in MA564 since many functions could be proven to be continuous by decomposing them componentwise into more obviously continuous maps, simplifying some tedious proofs. I would like to establish this result for smooth maps as well.

\begin{lemma}\label{3.2}
	There is a category $\mathcal{M} \mathrm{an} $ where objects are smooth manifolds and morphisms are smooth maps between manifolds.
\end{lemma}
\begin{proof}
	We proved in class that composites of smooth maps are smooth, so it suffices to show identities are smooth. This is obvious: given $1: M \to M$ with $M$ an $n$-manifold, let $p\in M$, and choose $(U,\varphi )$ a chart of $p$ in the relevant atlas. Of course since $1(p)=p$, the same chart works for the codomain. The composite $\varphi \circ 1\circ \varphi ^{-1} =\varphi \circ \varphi ^{-1} =1$ is smooth since $M$ is a manifold, and thus its charts are smoothly compatible. Additionally, $1(U)=U\subseteq U$.
\end{proof}

\begin{proposition}\label{3.4}
	$\mathcal{M} \mathrm{an} $ has finite products.
\end{proposition}
\begin{proof}
	We (probably, I don't remember; if not it's trivial) showed in class that the singleton $*$ admits a smooth structure, so it suffices to show $\mathcal{S} \mathrm{mth} $ has binary products. Let $M,N$ be $m$- and $n$-manifolds, respectively, and let $\mathcal{A} _M$, $\mathcal{A} _N$ be maximal atlases for $M$ and $N$, respectively. Define the product atlas
	\[
		\mathcal{A}_{M\times N} \coloneqq \left\{ (U\times V,\varphi \times \psi )\mid (U,\varphi)\in \mathcal{A} _M,(V,\psi )\in \mathcal{A} _V \right\} ,
	\]
	where the products (and all products in this problem, unless otherwise stated) are taken in $\Top $. Fairly obviously,
	\[
		\bigcup_{(U\times V,\varphi \times \psi )\in \mathcal{A} _{M\times N} } (U\times V)=\left( \bigcup_{(U,\varphi )\in\mathcal{A} _M} U \right) \times \left( \bigcup_{(V,\psi )\in \mathcal{A} _N} V \right) \supseteq M\times N,
	\]
	where the final step follows since $\mathcal{A}_M$ is an open cover of $M$, likewise for $N$, and simply by our knowledge of the cartesian product.

	It suffices to show the maps of charts are smoothly compatible. This indeed follows: Let $(U\times V,\varphi \times \psi ),(W\times X,\sigma \times \tau )\in \mathcal{A} _{M\times N} $. We show the composite
	\[\begin{tikzcd}[row sep = 2.5em, column sep = 4.5em]
	\widetilde{U} \times \widetilde{V} \ar[" (\varphi \times \psi )^{-1}  ", " \varphi ^{-1} \times \psi ^{-1}  "', r]  & U\times V \ar[" (\sigma \times \tau )|(U\times V) ", r]  & \widetilde{W} \times \widetilde{X} 
	\end{tikzcd}\]
	is continuous. The opposite direction is proven the same way. Of course if $ U \cap W$ or $V \cap X$ is empty, then the function is the empty function, so we assume without loss of generality that the relevant intersections are nonempty. Following a point along the composite gives
	\[\begin{tikzcd}[row sep = 2.5em, column sep = 2.5em]
		(p,q)\ar[mapsto, r]  & (\varphi ^{-1} (p),\psi ^{-1} (q)) \ar[mapsto, r]  & ((\sigma \varphi ^{-1}) (p),(\tau \psi ^{-1} )(q)).
	\end{tikzcd}\]
	Since $\sigma \varphi ^{-1} $ and $\tau \psi ^{-1} $ are smooth by hypothesis, it suffices to show the product of smooth functions between (open subsets of) $\mathbb{R} ^m\times \mathbb{R} ^n \to \mathbb{R} ^m \times \mathbb{R} ^n$ are smooth.

	To show this, we compute the partials of $(f\times g)$ for $f : \widetilde{U}  \to \widetilde{W} $ and $g : \widetilde{V}  \to \widetilde{X} $ smooth functions. Let $x\in \widetilde{U} $ and $y\in \widetilde{W} $, and write $x^i$, $y^j$ for the components of these vectors. Also write $f_i,g_j$ for the component functions of $f$ and $g$. Recall how partial derivatives work:
	\[
		\frac{\partial f}{\partial x^i}=\left( \frac{\partial f_1}{\partial x^i} , \ldots , \frac{\partial f_m}{\partial x^i}  \right) 
	.\]
	Of course, we then have
	\[
		\frac{\partial (f\times g)}{\partial x^i} =\left( \frac{\partial f_1}{\partial x^i} , \ldots ,\frac{\partial f_m}{\partial x^i} , \frac{\partial g_1}{\partial x^i} , \ldots , \frac{\partial g_n}{\partial x^i}  \right) =\left( \frac{\partial f_1}{\partial x^i} , \ldots , \frac{\partial f_m}{\partial x^i} , 0, \ldots ,0 \right) 
	\]
	where the 0s on the end are since $g$ does not depend on $x$. We obtain the obvious similar result when differentiating with respect to $y^j$. It should be clear at this point that since the constant map 0 is smooth, and since $f,g$ are smooth, we have that $f\times g$ is smooth. This suffices to prove that $\mathcal{A} _{M\times N} $ is an atlas for $M\times N$.
	
	Now we show $M\times N$ is a product. Let $f : Z \to M$ and $g : Z \to N$ be smooth maps. Define $(f,g): Z \to M\times N$ by $p\mapsto (f(p),g(p))$. This is the usual definition in $\Set $ so it is automatically universal in the relevant way, so it suffices to show $(f,g)$ and the projections are smooth with respect to the maximal atlas for $\mathcal{A} _{M\times N} $.

	We start by showing $(f,g)$ is smooth. Let $p\in Z$, and since $f,g$ are smooth, there are charts $(U,\varphi ),(V,\psi )$ for $f(p)$ and $g(p)$, respectively, and two charts $(W,\sigma ),(X,\tau )$ for $p\in Z$, such that
	\[
		f(W)\subseteq U,\qquad \varphi \circ f\circ \sigma^{-1}  \text{ smooth},
	\]
	\[
		g(X)\subseteq V,\qquad \psi \circ g\circ \tau^{-1}  \text{ smooth} 
	.\]
	Since $(W,\sigma ),(X,\tau )$ are smoothly related, we have that $\sigma \circ \tau ^{-1} ,\tau \circ \sigma ^{-1} $ are smooth. Since $p\in W \cap X$ is open, we have that $\psi \circ g\circ \tau ^{-1} \circ \tau \circ \sigma ^{-1} $ is smooth. Restricting along $W \cap X$, we thus obtain charts
	\[
		f(W \cap X)\subseteq U,\qquad \varphi \circ f \circ \sigma ^{-1} \text{ smooth},
	\]
	\[
		g(W \cap X)\subseteq V,\qquad \varphi \circ f\circ \sigma ^{-1} \text{ smooth}
	.\]
	The statement now follows fairly easily: $f(W \cap X)\times g(W \cap X)\subseteq U\times V$ obviously, and since $-\times -$ is a bifunctor on $\Set ,\Top $ with the same point-set construction of products, it preserves composition:
	\[
		(\varphi \circ f\circ \sigma ^{-1} )\times (\psi \circ g\circ \tau ^{-1} )=(\varphi \times \psi )\circ (f\times g)\circ (\sigma ^{-1} \circ \tau ^{-1} )
	.\]
	Since the composite on the left is smooth componentwise in the calc 3 sense, by our above proof it is smooth. By the equality of the two composites, the composite on the right is smooth, and thus we have a chart $(W \cap X \times W \cap X,\sigma \times \sigma )$ for $p$ and a chart $(U\times V,\varphi \times \psi )$ for $(f(p),g(p))$ witnessing $(f,g)$ as a smooth map.

	Finally, we show the projections are smooth. This is fairly direct. Let $(p,q)\in M\times N$, and consider the projection $\pi : M\times N \to M$. The proof for $N$ is the same. Choose a chart $(U\times V,\varphi \times \psi )$ for $(p,q)$, and since $p\in U$, we have a chart $(U,\varphi )$ for $p$ in $M$. Of course, $\pi (U\times V)=U\subseteq U$, so it suffices to prove $\varphi \circ \pi \circ (\varphi ^{-1} \times \psi ^{-1} )$ is smooth. This is fairly trivial: this is a map $\mathbb{R} ^{m+n} \to \mathbb{R} ^n$, so by our above proof it suffices to prove it is smooth componentwise. We have
	\[
		(\varphi (p),\psi (q))\mapsto (p,q)\mapsto p\mapsto \varphi (p)
	.\]
	This is clearly equivalent to the composite $\varphi \circ \varphi ^{-1} \circ \pi =\pi $:
	\[
		(\varphi (p),\psi (q))\mapsto \varphi (p)\mapsto p\mapsto \varphi (p)
	.\]
	Projections $\mathbb{R} ^{m+n} \to \mathbb{R} ^n$ are well-known to be smooth, so the projection $M\times N\to M$ is indeed smooth, completing the proof.
\end{proof}
\begin{corollary}\label{3.5}
	Product is a bifunctor $-\times - : \mathcal{M} \mathrm{an} \times \mathcal{M} \mathrm{an}  \to \mathcal{M} \mathrm{an} $.
\end{corollary}
\begin{proof}
	We specialize the usual proof to $\mathcal{M} \mathrm{an} $. Given $M_1,M_2,N_1,N_2$ smooth manifolds and smooth maps $f_1 : M_1 \to N_1$, $f_2 : M_2 \to N_2$, there is a diagram of solid arrows
	\[\begin{tikzcd}[row sep = 2.5em, column sep = 2.5em]
	M_1 \ar[" f_1 ", r]  & N_1 \\
	M_1\times M_2 \ar[u] \ar[" f_1\times f_2 ", dashed, r] \ar[d]  & N_1\times N_2\ar[u] \ar[d]  \\
	M_2\ar[" f_2 "', r]  & N_2
	\end{tikzcd}\]
	and by \cref{3.4}, there is a unique dashed arrow as indicated, rendering the diagram commutative. Preservation of 1 and $\circ $ follow manifestly from universality.
\end{proof}

The following are immediate by inspecting the proofs of \cref{3.4,3.5}.
\begin{corollary}\label{3.6}
	Let $f : M_1 \to N_1$ and $g : M_2 \to N_2$ be smooth maps. Then the map $f\times g : M_1\times M_2 \to N_1\times N_2$ given by $(m,n)\mapsto (f(m),g(n))$ is smooth.
\end{corollary}
\begin{corollary}\label{3.7}
	Let $f : M \to N_1$ and $g : M \to N_2$ be smooth maps. Then there is a smooth map $(f,g): M \to N_1\times N_2$ given by $m\mapsto (f(m),g(m))$.
\end{corollary}

We now return to the main result and drop the category-theoretics.

\begin{proof}[Proof of \cref{3}]
	(1) INCOMPLETE

	(2) Write $x=(x^1, \ldots ,x^{n+1} )$ and $x'=(x^1, \ldots ,x^n)$. Then
	\[
		(\sigma ^{-1} \circ \sigma )(x)=\sigma ^{-1} \left(\frac{x'}{1-x^{n+1} } \right)=\frac{(2x^1, \ldots , 2x^n, \left| x' \right| ^2-1}{1-x^{n+1} }\cdot \frac{1}{\left| x' \right| +1} 
	.\]
	INCOMPLETE

	(3) By \cref{3.4}, maps are smooth if and only if they are smooth componentwise. We have
	\[
		\widetilde{\sigma } \circ \sigma ^{-1} (x)=\widetilde{\sigma } \left( \frac{\left( 2x^1, \ldots ,2x^n, \left| x \right| ^2-1 \right) }{\left| x \right| ^2+1}  \right) =-\frac{1}{\left( \left| x \right| ^2+1 \right) \left( 1-\left(-\left| x \right| ^2-1\right) \right) } \left( -2x^1, \ldots ,-2x^n \right) 
	.\]
	Of course the result simlpifies to
	\[
		(\widetilde{\sigma } \circ \sigma ^{-1} )(x)=\frac{2x}{\left(2+\left| x \right| ^2\right)\left( \left| x \right| ^2+1 \right) } 
	.\]
	Note that addition, the norm, and multiplication are all smooth maps: the norm is the only one that might not be clear, but it is simply the square root of a sum of squares, all of which are clearly smooth. Composing smooth maps gives smooth maps, thus the entire composite is smooth. The transition map is thus componentwise smooth.

	(4) To show they define the same smooth structure, we show they are smoothly compatible. Since smooth manifolds take a maximal atlas, if they are smoothly compatible, then they always show up together in a maximal atlas. We show $\varphi ^{\pm} _j$ is smoothly compatible with $\sigma $; the proof for $\widetilde{\sigma } $ is similar. Recall $\varphi _i^{\pm} (x)=(x^1, \ldots ,x^{i-1} ,x^{i+1} ,\ldots x^{n+1})$. The composite $\varphi _i^{\pm} \circ \sigma ^{-1} $ is thus given by
	\[
		(\varphi _i^\pm\circ \sigma ^{-1} )(x)=\varphi _i^{\pm} \left( \frac{1}{\left| x \right| ^2+1}   \left( 2x^1, \ldots ,2x^n, \left| x \right| ^2-1 \right) \right)=\frac{1}{\left| x \right|^2 +1} \left( 2x^1, \ldots ,2x^{i-1} ,2x^{i+1} ,\ldots ,\left| x \right| ^2-1 \right) 
	.\]
	There are a few cases to check here. First for the case $i=n+1$, we obviously just have
	\[
		(\varphi ^{\pm} \circ \sigma ^{-1} )(x)=\frac{2x}{\left| x \right| ^2+1} 
	.\]
	This is clearly smooth componentwise for the same reasons as (3). Otherwise, all components except the final one are smooth by this same argument. For the final one, we note that norm and subtraction are smooth, showing that component is also smooth.

	For the other direction, we have
	\[
		(\sigma \circ \left( \varphi _i^{\pm}  \right) ^{-1} )(x)=\sigma (x^1, \ldots ,x^{i-1} ,1-\left| x \right| ,x^{i+1} ,\ldots ,x^n)=\frac{1}{1-x^n} (x^1, \ldots ,x^n),
	\]
	where we have taken $i<n$. This is smooth by the same arguments as before. If $i=n$, then we instead have
	\[
		(\sigma \circ \left( \varphi _i^{\pm}  \right) ^{-1} )(x)=\frac{1}{\left| x \right| } (x^1, \ldots ,x^{n-1} , 1-\left| x \right| ),
	\]
	which again is smooth by the same logic.
\end{proof}


\section{Problem 1.9}

\begin{problem}\label{4}
	Show that $\mathbb{C} \mathbb{P} ^n$ is a compact $2n$-manifold, and has a smooth structure analogous to $\mathbb{R} \mathbb{P} ^n$. We identify $\mathbb{C} ^{n+1} $ with $\mathbb{R} ^{2n+2} $ by $(x^1+iy^1, \ldots ,x^{n+1} +iy^{n+1} )\mapsto (x^1,y^1, \ldots ,x^{n+1} ,y^{n+1} )$.
\end{problem}
\begin{proof}
	The same proof for $\mathbb{R} \mathbb{P} ^n$ will work for showing that $\mathbb{C} \mathbb{P} ^n$ is a $2n$-manifold, just with the extra consideration of the homeomorphism $\mathbb{C} ^n\cong \mathbb{R} ^{2n} $. Define
	\[
		\widetilde{U} _i\coloneqq \left\{ x\in\mathbb{C} ^{n+1} \mid x^i\neq 0 \right\}\subseteq \mathbb{C} ^{n+1} \setminus \left\{ 0 \right\}  ,\qquad U_i\coloneqq \pi (\widetilde{U} _i)
	\]
	for $\pi : \mathbb{C} ^{n+1} \setminus \left\{ 0 \right\} \to \mathbb{C} \mathbb{P} ^n$ the quotient map. Since $x^i\neq 0$, we may consider a map $\varphi $
	\[
		[x]\mapsto \left( \frac{x^1}{x^i} , \ldots ,\frac{x^{i-1} }{x^i} ,\frac{x^{i+1} }{x^i} ,\ldots ,\frac{x^n}{x^i} \right)\in\mathbb{C} ^{n} .
	\]
	Of course, this map may be composed with the homeomorphism $\mathbb{C} ^n \cong \mathbb{R}^{2n} $ noted in the problem, and since composites of invertible maps are invertible, it suffices to show the given map is a homeomorphism $U_i\cong \mathbb{C} ^n$. This involves first showing the map is well-defined.

	Let $y\sim x\in U_i$ in $\mathbb{C} ^{n+1} \setminus \left\{ 0 \right\} $, meaning $x^i=\lambda y^i$ for some fixed nonzero $\lambda \in\mathbb{C} $. We then have
	\[
		\varphi [y]=\left( \frac{\lambda x^1}{\lambda x^i} , \ldots , \frac{\lambda x^{i-1} }{\lambda x^i} ,\frac{\lambda x^{i+1} }{\lambda x^i} ,\ldots ,\frac{\lambda x^n}{\lambda x^i}  \right) =\varphi [x]
	.\]
	Therefore, the map is well-defined. We define an inverse map $\varphi ^{-1} : \mathbb{C} ^n \to U_i$ by
	\[
		x\mapsto \left[ x^1, \ldots ,x^{i-1} ,1,x^{i+1} ,\ldots ,x^n \right] 
	.\]
	Of course, we have that for any $[x]\in U_i$,
	\[
		\varphi ^{-1} \varphi [x]=\left[ \frac{x^1}{x^i} , \ldots ,\frac{x^{i-1} }{x^i} ,1,\frac{x^{i+1} }{x^i} ,\ldots ,\frac{x^n}{x^i}  \right] =\left[ \frac{1}{x^i} x \right] =[x]
	\]
	since $\frac{1}{x^i} x\sim x$ by definition. For the other side, we have for any $x\in\mathbb{C} ^n$,
	\[
		\varphi \varphi ^{-1} (x)=\varphi \left[ x^1, \ldots ,x^{i-1} , 1, x^i, \ldots ,x^n \right] =\left( x^1, \ldots ,x^n \right) =x
	.\]
	It suffices at this point to show $\varphi ,\varphi ^{-1} $ are continuous. By Lee A.23 (a), maps to product spaces are continuous if and only if they are continuous componentwise. Likewise, by Lee A.27 (a), maps \emph{from} quotient spaces are continuous if and only if the corresponding map from the unquotiented space is continuous. In less words, Lee describes universality of products and quotients, which I use to make continuity of these maps manifestly obvious.

	By these theorems, $\varphi $ is continuous if and only if the corresponding map $\mathbb{C} ^{n+1}\setminus \left\{ 0 \right\} \to \mathbb{C} ^n$ is continuous, which is continuous if and only if each component function is continuous\footnote{Of course $\mathbb{C} ^{n+1} \setminus \left\{ 0 \right\} $ is not a product space since it is the subspace of a product space, but we may view $\varphi $ as a map $\mathbb{C} ^{n+1} \to \mathbb{C} $ by defining $0\mapsto 0$, then composing along the inclusion $\mathbb{C} ^{n+1} \setminus \left\{ 0 \right\} \hookrightarrow \mathbb{C} ^{n+1} $ to obtain the claimed result.}. Each component function can be described as the composite of a projection onto a component of the product, and division by a fixed number $x^i$, both of which are well-known to be continuous.

	Conversely, we decompose $\varphi ^{-1} $ as $\pi \circ \psi $, where
	\[
		\psi : x \mapsto (x^1, \ldots ,x^{i-1} ,1,x^{i+1} ,\ldots ,x^n)
	.\]
	Since $\pi $ is continuous, it suffices to show $\psi $ is continuous, which follows if and only if it is continuous componentwise. Indeed, the identity is manifestly continuous, as is the constant map at 1. Finally, note also that since the image under $\psi $ has $x^i\neq 0$ in $\mathbb{C} ^{n+1}\setminus \left\{ 0 \right\} $, the image under $\psi $ is necessarily an element of $\widetilde{U} _i$, and thus descends to $U_i$ under $\pi $.

	Switching notation: write $\varphi _i$ for the map $U_i \cong \mathbb{C} ^n$ composed with the homeomorphism $\mathbb{C} ^n \cong \mathbb{R} ^{2n} $. For $(U_i,\varphi _i)$ to be a chart, we need $U_i$ to be open. Note that the preimage of $U_i$ is at least $\widetilde{U} _i$. Moreover, if $p\in U_i$, then it is the image of some $\widetilde{p} \in \widetilde{U} _i$, so it suffices to show for any $\widetilde{p} \in \widetilde{U} _i$, the entire linear subspace $\operatorname{span} \widetilde{p} $ spanned by that point is a subset of $\widetilde{U} _i$. Indeed, if $\widetilde{p} \in U_i$, then $x^i\neq 0$, and thus $\lambda x^i\neq 0$ for all $\lambda \in\mathbb{C} \setminus \left\{ 0 \right\} $. Therefore, the preimage of $U_i$ is precisely $\widetilde{U} _i$. Note that the distance from any point $\widetilde{p} \in \widetilde{U} _i$ to the $x^i$-axis is simply $\pi _i(\widetilde{p} )$ for $\pi _i$ the projection onto the $i$th component of the product. We thus clearly have an open cover of $\widetilde{U} _i$ by open balls
	\[
		\widetilde{U} _i = \bigcup_{\widetilde{p} \in \widetilde{U}_i} B_{\pi _i(\widetilde{p} )} (\widetilde{p} )
	.\]
	Therefore, $\widetilde{U} _i$ is open, and $(U_i,\varphi _i)$ is a chart, and $\mathbb{C} \mathbb{P} ^n$ is a $2n$-manifold. We now just need to show each $\varphi _i$ is smoothly compatible, and that the $U_i$s form an open cover of $\mathbb{C} \mathbb{P} ^n$. The latter is fairly simple: given some $[x]\in\mathbb{C} \mathbb{P} ^n$, there is at least one $x^i\neq 0$ in the representative $x\in\mathbb{C} ^{n+1} \setminus \left\{ 0 \right\} $, and thus $[x]\in U_i$. For the latter, let $x\in\mathbb{R} ^{2n} $, and follow $x$ along the composite $\varphi _j\circ \varphi _i^{-1} $:
	\begin{align*}
		x&\mapsto \left(x^1+ix^2, x^3+ix^4, \ldots ,x^{2n-1} +ix^{2n} \right)\\
		 &\mapsto \left[ x^1+ix^2, \ldots ,x^{2i-2} +ix^{2i-1} ,1,x^{2i} +ix^{2i+1} , \ldots ,x^{2n-1} +ix^{2n}  \right] \\
		 &=\left( x^{2j} +ix^{2j+1}  \right) \left[ x^1+ix^2, \ldots ,x^{2i-2} +ix^{2i-1} ,1,x^{2i} +ix^{2i+1} , \ldots ,x^{2n-1} +ix^{2n}  \right] \\
		 &\mapsto \left( x^1+ix^2, \ldots ,x^{2i-2} +ix^{2i-1} ,x^{2j} +ix^{2j+1} ,x^{2i} +ix^{2i+1} ,\ldots ,x^{2j-2} +ix^{2j-1} ,x^{2j+2} +ix^{2j+3},\right.\\
		 &\qquad\left.\ldots ,x^{2n-1}+ix^{2n}   \right) \\
		 &\mapsto \left( x^1, \ldots ,x^{2i-1} ,x^{2j} ,x^{2j+1} ,x^{2i} ,\ldots ,x^{2j-1} ,x^{2j+2} ,\ldots ,x^{2n}  \right) 
	.\end{align*}
	It should be fairly clear that this map is smooth, since any partial $\partial _k$ will simply map all entries to 0, except the entry that $x^k$ ends up in, which will be 1. Constant maps are smooth, and thus the entire map is smooth, since all it does is swap around where a few entries live. Therefore, there is an atlas
	\[
		\mathcal{A} =\left\{ (U_i,\varphi _i) \right\}_{0\leq i\leq 2n} 
	.\]
	Therefore, $\mathbb{C} \mathbb{P} ^n$ admits a smooth structure.

	Realizing now I forgot to show $\mathbb{C}\mathbb{P} ^n$ is Hausdorff and second countable. We show these now. Let $[x]\neq [y]\in\mathbb{C} \mathbb{P} ^n$, and choose representatives $x\neq y\in\mathbb{C} ^{n+1} $. Since $[x]\neq [y]$, $x$ and $y$ are linearly independent. I'm running out of time to submit this assignment so I informally describe how the Hausdorff proof follows. $\left\{ x,y \right\} $ form a basis for a 2-dimensional subsace of $\mathbb{C} ^{n+1} $, so we normalize them in this basis and find the angle between the subspaces they define. We find the maximum radii of $x$ and $y$ such that balls of this radii do not cross this angle, then we quotient down to $\mathbb{C} \mathbb{P} ^n$. Since these balls did not cross the middle angle, the spans of each individual vector are all disjoint, meaning we have constructed disjoint sets containing $[x]$ and $[y]$. We now need to show they are open, which we do by the same exact process done below in the proof of second-countability.

	To show $\mathbb{C} \mathbb{P} ^{n} $ is second-countable, we note that $\mathbb{C} ^{n+1} \cong \mathbb{R} ^{2n+2} $ is second-countable, and thus by appealing to \cref{1.4} it suffices to show $\pi $ is open. Let $B_r(x)$ be an open ball in $\mathbb{C} ^{n+1} $, and note that
	\[
		\pi ^{-1} \pi (B_r(x))=\left\{ \lambda z\in\mathbb{C} ^{n+1} \mid z\in B_r(x),\lambda \in\mathbb{C} \setminus \left\{ 0 \right\}  \right\} 
	.\]
	Indeed, note that $\lambda B_r(x)=\left\{ \lambda z\mid z\in B_r(x) \right\} $ is open for $\lambda \in\mathbb{C} \setminus \left\{ 0 \right\} $, and by definition $\lambda B_r(x)\subseteq \pi ^{-1} \pi (B_r(x))$. This is also open since multiplication by a fixed nonzero constant is a homeomorphism. We then clearly obtain
	\[
		\pi ^{-1} \pi (B_r(x))=\bigcup_{\lambda \in\mathbb{C} \setminus \left\{ 0 \right\} } \lambda B_r(x),
	\]
	and thus $\pi ^{-1} \pi (B_r(x))$ is open. Thus, $\mathbb{C} \mathbb{P} ^n$ is second countable.
\end{proof}

\section{Problem 2.3}
\begin{problem}\label{5}
	Compute sufficiently many coordinate representations to prove the following maps are smooth.
	\begin{enumerate}
		\item $p_n : S^1 \to S^1$ for each $n\in\mathbb{Z} $ given by $z\mapsto z^n$, viewing $S^1 = e^{i \pi [0,1]}  \subseteq \mathbb{C} $.
		\item $\alpha : S^n \to S^n$ the antipodal map $x\mapsto -x$.
		\item $f : S^3 \to S^2$ where
			\[
				S^3 = \left\{ (w,z)\in\mathbb{C} ^2\mid \left| w \right| ^2+\left| z \right| ^2=1 \right\} \subseteq \mathbb{C} ^2,
			\]
			\[
				f(w,z) = (z \overline{w} +w \overline{z} ,iw \overline{z} -iz \overline{w} , z \overline{z} -w \overline{w} )
			.\]
	\end{enumerate}
\end{problem}

\begin{proof}
	(1) The atlas is generated by $\varphi ^+_1(x^1,x^2)=x^1$, and likewise for $\varphi _1^-,\varphi _2^{\pm} $. Choose $p=e^{i\theta } \in S^1$, and without loss of generality choose it in the upper hemisphere. Then $\varphi _1^{+,-1} (\cos \theta )=p$ kj sndkjdjksn INCOMPLETE
\end{proof}


\section{Problem 2.6}

\begin{problem}\label{6}
	Let $p : \mathbb{R} ^{n+1} \setminus \left\{ 0 \right\}  \to \mathbb{R} ^{k+1} \setminus \left\{ 0 \right\} $ be smooth, and suppose for some $d\in\mathbb{Z} $, $p(\lambda x)=\lambda ^dp(x)$ for all $\lambda \in\mathbb{R} \setminus \left\{ 0 \right\} $ and $x\in\mathbb{R} ^{n+1} \setminus \left\{ 0 \right\} $. Show the induced map $\widetilde{p} : \mathbb{R} \mathbb{P} ^n \to \mathbb{R} \mathbb{P} ^k$ is well-defined and smooth.
\end{problem}
\begin{proof}
	Lee A.27 describes universality of quotient in $\Top $. By this property, if we can show $p$ preserves equivalence classes $\mathbb{R} ^{n+1} \to \mathbb{R} ^{k+1} $, then it will descend to a map
	\[\begin{tikzcd}[row sep = 2.5em, column sep = 2.5em]
	\mathbb{R} ^{n+1} \ar[" p ", r] \ar[" \pi "', d]  & \mathbb{R} ^{k+1} \ar[" \pi  ", d]  \\
	\mathbb{R} \mathbb{P}^{n} \ar[" \widetilde{p}  "', r]  & \mathbb{R} \mathbb{P}^{k} 
	\end{tikzcd}\]
	We start by showing this. We must show if $x=\lambda x$, then $p(x)=k p(x)$ for some $\lambda ,k\in\mathbb{R} $. This follows by definition:
	\[
		p(\lambda x)=\lambda ^dp(x),\qquad \lambda ^d\in\mathbb{R} 
	.\]
	Thus, the map descends as needed, and $\widetilde{p} $ is well-defined.

	Now we show $\widetilde{p} $ is smooth. Let $(U_i,\varphi _i)$ be a chart for $x$, and $(V_j,\psi _j)$ a chart for $\widetilde{p} [x]$. By our explicit description of $\varphi _i$, we have
	\[
		(\psi \circ \widetilde{p} \circ \varphi _i ^{-1} )(x) = (\psi \circ \widetilde{p} )[x^1, \ldots ,x^{i-1} ,1,x^{i+1} , \ldots ,x^{n+1} ]=\psi \left[ p_1\varphi _i^{-1} (x), \ldots , p_{n+1} \varphi _i^{-1} (x) \right] 
	\]
	\[
	=\left( \frac{p_1 \varphi_i ^{-1} (x)}{p_j\varphi _i^{-1} (x)} , \ldots , \frac{p_{j-1} \varphi ^{-1} _i(x)}{ p_j\varphi _i^{-1} (x)}, \ldots , \frac{p_{j+1} \varphi ^{-1} _i(x)}{ p_{j+1}\varphi _i^{-1} (x)}, \ldots ,\frac{p_{n+1} \varphi ^{-1} _i(x)}{ p_{n+1}\varphi _i^{-1} (x)}\right) 
	.\]
	Note that each $p_\mu $ is smooth since a map is smooth if and only if it is smooth componentwise. Similarly, $\varphi _i^{-1} $ and $\psi $ (more precisely, division) are all already known to be smooth, so each component is a composite of smooth functions, and thus $\psi_j \circ \widetilde{p} \circ \varphi _i ^{-1} $ is smooth. Thus, $\widetilde{p} $ is smooth.
\end{proof}

\section{Problem 2.9}
\begin{problem}\label{7}
	Given a nonzero polynomial $p$ in one variable with complex coefficients, show there is a unique smooth map $\widetilde{p} : \mathbb{C} \mathbb{P} ^1 \to \mathbb{C} \mathbb{P} ^1$ as indicated by the dashed arrow
	\[\begin{tikzcd}[row sep = 2.5em, column sep = 2.5em]
	\mathbb{C} \ar[" G ", r] \ar[" p "', d]  & \mathbb{C} \mathbb{P}^1 \ar[" \widetilde{p}  ", dashed, d]  \\
	\mathbb{C} \ar[" G "', r]  & \mathbb{C} \mathbb{P} ^1
	\end{tikzcd}\]
	rendering the diagram commutative, where $G(z)=[z,1]$.
\end{problem}
\begin{proof}
	I assume we're allowed to assume $G$ is smooth. Given $z\in\mathbb{C} $, we have $Gp(z)=[p(z),1]$, and thus we are forced to define $\widetilde{p} $ such that $[z,1]\mapsto [p(z),1]$. This does not suffice as a full definition for $\widetilde{p}$, which we now explore.

	Let $[x,y]\in\mathbb{C} \mathbb{P} ^1$ with $y\neq 0$. Then $[x/y, 1]=[x,y]$ since it is given by scalar multiplication by $1/y$, so this formula defines $\widetilde{p} $ on all classes \emph{except} the class of $[1,0]$. The diagram commutes no matter what this definition is, since restricting along $G$ never has $[1,0]$ in the domain of $\widetilde{p} $:
	\[\begin{tikzcd}[row sep = 2.5em, column sep = 2.5em]
		z \ar[r, mapsto] \ar[d, mapsto]  & {[z,1]}\ar[d, mapsto]  \\
		p(z)\ar[r, mapsto]  & {[p(z),1]}
	\end{tikzcd}\]
	We will now derive what $\widetilde{p} [1,0]$ should be to complete the definition. Since $p$ is a complex polynomial, write $p=a_\mu z^\mu $ (somewhat using Einstein notation because I'm lazy), with maximal degree $d$. We have
	\[
		[p(x/y),1]=\left[ a_\mu \frac{x^\mu }{y^\mu } ,1 \right] =\left[ y^d \sum_{i=0}^{d} a_i \frac{x^i}{y^i} ,y^d \right] =\left[ \sum_{i=0}^{d} a_ix^iy^{d-i} ,y^d \right] 
	.\]
	Writing $P(x,y)=a_\mu x^\mu y^{d-\mu } $, we have
	\[
		\widetilde{p} [x,y]=\widetilde{p} [x/y,1]=[p(x/y),1]=[P(x,y),y^d]
	.\]
	This formulation allows us to extend $\widetilde{p} $ to $[1,0]$:
	\[
		\widetilde{p} [1,0]=[P(1,0),0]=\left[ a_d,0 \right] =[1,0]
	.\]
	
	Now we show this definition of $\widetilde{p} $ is smooth. Write $U_1=\left\{ [x,y]\mid x\neq 0 \right\} $, and $U_2$ for the other chart. For any $[z,1]\in U_2$, we have the chart $(U_2,\varphi _2)$. Note that $(U_2,\varphi _2)$ is also a chart for $\widetilde{p} [z,1]=[p(z),1]$, and we also immediately see that $\widetilde{p} (U_2)\subseteq U_2$. Consider the composite $\varphi _2 \circ \widetilde{p} \circ \varphi _2^{-1} $
	\[
		z\mapsto [z,1]\mapsto [p(z),1]\mapsto p(z).
	\]
	The map $z\mapsto p(z)$ is smooth because $p$ is smooth.

	Now consider the point $[1,0]\in U_1$, which has a chart $(U_1,\varphi _1)$. Since $\widetilde{p} [1,0]=[1,0]$, the same chart works for the image. Consider the composite $\varphi _1 \circ \widetilde{p} \circ \varphi _1^{-1} $:
	\[
		z\mapsto [1,z]\mapsto [P(1,z),z^d]\mapsto z^d,
	\]
	which is smooth because exponentials are smooth. I don't know why $\widetilde{p} (U_1)\subseteq U_1$. Therefore, $\widetilde{p} $ is smooth. Uniqueness is once again forced by the commutativity condition and the extension function $P(-,-)$.
\end{proof}

\section{Problem 2.10}

\begin{problem}\label{8}
	Let $M$ be a topological space, and let $C(M)$ be the algebra of continuous maps $M\to \mathbb{R} $. Given a continuous map $F : M \to N$, let $F^* : C(N) \to C(M)$ denote precomposition.
	\begin{enumerate}
		\item Show that $F^*$ is linear.
		\item Let $M,N$ be smooth manifolds. Show that $F : M \to N$ is smooth if and only if $F^*(C^\infty (N))\subseteq C^\infty (M)$.
		\item Let $F$ be a homeomorphism $M\cong N$. Show $F$ is a diffeomorphism if and only if $F^*: C^\infty (N) \to C^\infty (M)$ is an isomorphism.
	\end{enumerate}
\end{problem}
\begin{proof}
	(1) Let $f : N \to \mathbb{R} $, and recall that $cf$ for $c\in\mathbb{R} $ is defined by $(cf)(x)=c\cdot f(x)$. We have
	\[
		F^*(cf)(x)=(cf\circ F)(x)=(cf)(F(x))=c\cdot f(F(x))=c\cdot (f\circ F)(x)=cF^*(f)(x)
	.\]
	Now let $f,g : N \to \mathbb{R} $, and recall that addition is defined pointwise $(f+g)(x)=f(x)+g(x)$. We then have
	\[
		F^*(f+g)(x)=((f+g)\circ F)(x)=(f+g)(F(x))=f(F(x))+g(F(x))=(f\circ F)(x)+(g\circ F)(x)
	\]
	\[
		=F^*(f)(x)+F^*(g)(x)
	.\]
	Therefore, $F^*$ is $\mathbb{R} $-linear.

	(2) Let $F : M \to N$ be smooth. Then if $f : N \to \mathbb{R} $ is smooth, then since composites of smooth maps are smooth, $F^*(f)=f\circ F : M \to \mathbb{R} $ is smooth, so $F^*(C^\infty (N))\subseteq C^\infty (M)$. For the converse, let $F^*(C^\infty (N))\subseteq C^\infty (M)$. Let $p\in M$, and let $(U,\varphi )$ be a chart for $p$ and $(V,\psi )$ a chart for $F(p)$. Note that $\varphi ^{-1} $ is smooth $\hat{U} \to U$. If we can show that $\psi $ extends to a smooth map $\psi ' : N\to \mathbb{R} ^n$, we can postcompose by a projection to $\mathbb{R} $ which is aso smooth, and thus the composite $\pi \circ \psi '\circ F\circ \varphi ^{-1} $ will be a composition of smooth functionns, since $\pi \circ \psi '\circ F=F^*(\pi \circ \psi ')$ and $\pi \circ \psi '$ is smooth. It suffices to show this. I did not show this.

	(3) This follows by (2). Since $F$ is a homeomorphism, it has an inverse homeomorphism $F^{-1} : N \to M$. If $F$ is a diffeomorphism, then $F,F^{-1} $ are smooth. By (2), it follows that $F^*(C^\infty (N))\subseteq C^\infty (M)$, and $(F^{-1} )^*(C^\infty (M))\subseteq C^\infty (N)$. The maps $F^* : C^\infty (N) \to C^\infty (M)$, $(F^{-1} )^*: C^\infty (M) \to C^\infty (N)$ are thus well-defined, and the functions are inverses:
	\[
		(F^{-1} )^*F^*(f)=f\circ F\circ F^{-1} =f,\qquad F^*(F^{-1} )^*(f)=f\circ F^{-1} \circ F=f
	.\]
	Thus, $F^*$ is an isomorphism. For the converse direction, let $F^*$ be an isomorphism $C^\infty (N)\to C^\infty (M)$, and thus the only possible inverse is $(F^{-1} )^*$. We thus have that $F^*(C^\infty (N))=C^\infty (M)\subseteq C^\infty (M)$, and similarly for $F^{-1} $, and thus $F,F^{-1} $ are smooth, and thus diffeomorphisms.
\end{proof}

\end{document}
